% section-07.tex
% Corrected Section 10.7 (Feynman symbolic regression: PCA-split discussion)
% Written for use with jmlr2e.sty (loads epsfig, amssymb, natbib, graphicx).
%
% NOTE: jmlr2e.sty does NOT load the todonotes package, so \todo{} is not
% available unless you add \usepackage{todonotes} yourself in the main
% document preamble. To avoid a compile failure, the DEFAULT block below
% uses plain LaTeX comments instead of \todo{}. A todonotes version is
% included further down as an opt-in alternate -- only use it if
% \usepackage{todonotes} has been added to your main .tex file.
%
% CONTENT NOTE: The original draft claimed a "reduction" from 0.678 to
% 0.789, but 0.789 > 0.678 (an increase), and 0.678 does not match any
% figure derivable elsewhere in the paper (9/30 = 0.300, matching
% abstract/Table 7, or 142/180 = 0.789, the PCA-split corrected result).
% Two items below still require author verification before this is
% submission-ready:
%   1. INSERT_BASELINE / INSERT_SOURCE - the true, verified random-split
%      baseline number and where it comes from.
%   2. The explanation of why a harder split yields a higher score, and
%      the relationship between this evaluation and the primary result.

\subsection{Feynman Symbolic Regression: PCA-Split Correction}
\label{sec:feynman-pca-correction}

The increase from the random-split baseline
($\mathrm{INSERT\_BASELINE}$)  % TODO: verify value + cite source (INSERT_SOURCE)
to the PCA-directed split result ($142/180 = 0.789$) reflects
% TODO: explain why a harder split yields a numerically higher score
% here -- e.g., a different metric, a different evaluation subset, or
% a correction to an earlier miscalculation.
\ldots
This is distinct from the primary result reported in the abstract
and Table~\ref{tab:primary-results} ($9/30 = 0.300$), which
% TODO: clarify the relationship between the two evaluations if they
% are not directly comparable.
\ldots

% -----------------------------------------------------------------
% OPT-IN ALTERNATE: todonotes version (visible TODO markers in the
% typeset PDF). Only use this block if \usepackage{todonotes} has
% been added to your main document's preamble; otherwise it will
% fail to compile. To use: delete the \subsection block above and
% uncomment the block below.
% -----------------------------------------------------------------
%
% \subsection{Feynman Symbolic Regression: PCA-Split Correction}
% \label{sec:feynman-pca-correction}
%
% The increase from the random-split baseline
% ($\mathrm{INSERT\_BASELINE}$\todo{Verify: source = INSERT\_SOURCE})
% to the PCA-directed split result ($142/180 = 0.789$) reflects
% \todo{Explain why a harder split still yields a numerically higher
% score here --- e.g., a different metric, a different evaluation
% subset, or a correction to an earlier miscalculation.}
% This is distinct from the primary result reported in the abstract
% and Table~\ref{tab:primary-results} ($9/30 = 0.300$), which
% \todo{Clarify the relationship between the two evaluations if they
% are not directly comparable.}
